\section*{Chapitre \ref{ch:nw:disaster-recovery-and-regulatory-requirements} --- Reprise après sinistre et exigences réglementaires}

\begin{solution}{exo:nw:disaster-recovery-and-regulatory-requirements:1}
Reprise du trading le jour ouvré suivant et reprise en deux heures des systèmes SCI critiques à la suite d'une perturbation de grande ampleur,
avec des capacités de secours suffisamment résilientes et géographiquement diversifiées.
\end{solution}

\begin{solution}{exo:nw:disaster-recovery-and-regulatory-requirements:2}
L'\omterm{def:nw:disaster-recovery-and-regulatory-requirements:objectives}{objectif de délai de reprise} (RTO) borne la durée pendant laquelle un système peut rester à l'arrêt~; l'\omterm{def:nw:disaster-recovery-and-regulatory-requirements:objectives}{objectif de point de reprise} (RPO)
borne la quantité de données récentes qui peut être perdue. Le premier est fixé par le runbook, le second par la réplication.
\end{solution}

\begin{solution}{exo:nw:disaster-recovery-and-regulatory-requirements:3}
Il se trouve à \qty{34.2}{\kilo\metre} de Secaucus, dans le périmètre de toute catastrophe régionale qui toucherait le site principal~:
une tempête, une inondation ou une panne de courant toucheraient probablement les deux.
\end{solution}

\begin{solution}{exo:nw:disaster-recovery-and-regulatory-requirements:4}
La réplication ne fixe que le point de reprise~; le temps du dispositif warm est celui de son runbook, dominé par le démarrage et la restauration des serveurs ainsi que par le chargement des données,
lesquels ne sont pas modifiés par une réplication plus rapide.
\end{solution}

\begin{solution}{exo:nw:disaster-recovery-and-regulatory-requirements:5}
Dispositions de remplacement pour gérer les ordres et positions en cours, et dispositions permettant d'arrêter les algorithmes ou les systèmes
sans créer de conditions de négociation désordonnées.
\end{solution}

\begin{solution}{exo:nw:disaster-recovery-and-regulatory-requirements:6}
Oui pour le 28 septembre 2026~; non pour le 1\up{er} novembre 2026, plus de douze mois après le test.
\end{solution}

\begin{solution}{exo:nw:disaster-recovery-and-regulatory-requirements:7}
De 20.3~\% à 44.1~\%~: l'étape de restauration est la plus longue, mais les cinq autres étapes, prises ensemble à leur médiane, prennent encore environ deux heures.
\end{solution}

\begin{solution}{exo:nw:disaster-recovery-and-regulatory-requirements:8}
Un test n'est pas une décision~: avant Sandy, le test du secteur avait réussi deux jours plus tôt et les marchés n'ont pourtant pas basculé.
L'entreprise a besoin d'un runbook doté d'un responsable de la décision et de critères, d'une bascule répétée, et non seulement de la connectivité, ainsi que des plans propres aux bourses.
\end{solution}

\begin{solution}{pb:nw:disaster-recovery-and-regulatory-requirements:1}
\textbf{Partie~I.}
\begin{enumerate}
\item NY4 \qty{0.3}{\kilo\metre}, Carteret 26.0, Mahwah 34.2, Aurora 1\,191.1~; seul Aurora est hors du périmètre.
\item \qty{17.4}{\milli\second} d'aller-retour vers Aurora avec un \omterm{def:nw:the-north-american-map:factor}{facteur de tracé} de 1.5, \qty{0.5}{\milli\second} vers Mahwah.
\item Être hors des catastrophes susceptibles de frapper le site principal, les règles exigeant une diversité géographique.
\item \qty{56}{\milli\second}~: \qty{50}{\milli\second} de retard, auquel s'ajoute le délai à sens unique.
\end{enumerate}
\textbf{Partie~II.}
\begin{enumerate}[resume]
\item 99.2~\%, 64 et 96 minutes.
\item 20.3~\%, 148 et 237 minutes.
\item Au plus quinze minutes, l'intervalle entre deux instantanés.
\item Démarrage et restauration des serveurs (45 minutes à la médiane) et chargement des données (30).
\end{enumerate}
\textbf{Partie~III.}
\begin{enumerate}[resume]
\item Des plans raisonnablement conçus pour reprendre les systèmes critiques dans les deux heures et le trading le jour ouvré suivant, géographiquement diversifiés~;
  des membres désignés testant avec la bourse au moins tous les douze mois.
\item Dispositions de continuité d'activité documentées par plateforme, délocalisation vers un site de secours le cas échéant, arrêt ordonné, gestion
  des ordres et positions en cours, revue annuelle et test.
\item Les entreprises devaient pouvoir rester dans les \omterm{def:nw:disaster-recovery-and-regulatory-requirements:resilience}{tolérances d'impact} de leurs \omterm{def:nw:disaster-recovery-and-regulatory-requirements:resilience}{services métier importants}.
\item Le risque que les entreprises et les bourses ne basculent pas, même lorsque les sauvegardes fonctionnaient~: la décision et la confiance comptent
  autant que le matériel.
\end{enumerate}
\textbf{Partie~IV.}
\begin{enumerate}[resume]
\item \emph{Résultat nommé}~: avec les runbooks du chapitre, le dispositif hot reprend en deux heures dans 99.2~\% des cas et le dispositif warm
  dans 20.3~\%~; à 70~\% et 25~\% de 3.19~millions de dollars par an du chapitre~27, l'écart coûte 1.44~millions de dollars par an.
\item Publié~: les règles, les sites de reprise des bourses, les dates de test et les coordonnées des sites. Hypothèses~: la durée de chaque étape, les coûts,
  le rayon du danger, le décalage de réplication.
\item Si la \omterm{def:nw:disaster-recovery-and-regulatory-requirements:resilience}{tolérance d'impact} de l'entreprise pour le trading est de deux heures, et qu'un jour perdu coûte souvent plus de 1.44~millions de dollars~;
  sinon, un site warm avec arrêt ordonné peut être le choix proportionné.
\item Serveurs préparés à l'avance, restaurations automatisées, étapes en parallèle, et runbooks mesurés et répétés.
\item Services et tolérances nommés, runbook, résultats de tests avec des durées mesurées, et traces de la participation aux tests des bourses.
\item En tant que membre désigné lorsqu'on le demande, et sinon de manière volontaire, avec le même runbook que pour une situation réelle.
\item Dans les lignes de résilience~: le site de reprise, ses lignes et ses \omterm{def:nw:colocation-products-and-how-they-are-sold:mmr}{interconnexions} (cross-connects), la réplication, le runbook et le calendrier des tests.
\item Heures de trading récupérées après un sinistre, à un prix annuel.
\end{enumerate}
\end{solution}

\begin{solution}{iq:nw:disaster-recovery-and-regulatory-requirements:1}
L'\omterm{def:nw:disaster-recovery-and-regulatory-requirements:objectives}{objectif de délai de reprise} (RTO) est la durée pendant laquelle un système peut rester à l'arrêt~; l'\omterm{def:nw:disaster-recovery-and-regulatory-requirements:objectives}{objectif de point de reprise} (RPO)
est la quantité de données qui peut être perdue. Le runbook fixe le premier~; la réplication fixe le second.
\iqlookfor{Définitions~; runbook face à la réplication.}
\end{solution}

\begin{solution}{iq:nw:disaster-recovery-and-regulatory-requirements:2}
Hors des risques partagés de la zone métropolitaine de New York, à proximité des sites de reprise propres aux bourses et bien reliés à eux~: Aurora
est le choix naturel pour les contrats à terme américains et BrokerTec~; le coût tient à la réplication asynchrone et à un point de reprise.
\iqlookfor{Risques~; sites de reprise des plateformes~; compromis de réplication.}
\end{solution}

\begin{solution}{iq:nw:disaster-recovery-and-regulatory-requirements:3}
Une seule autorité décide quel site est actif (une personne munie de critères, ou un quorum), l'ancien principal est isolé avant que le nouveau
ne négocie, et les sessions utilisent les adresses de reprise distinctes de la bourse, de sorte que les deux sites ne puissent pas négocier simultanément.
\iqlookfor{Décision unique~; isolement~; adresses distinctes.}
\end{solution}

\begin{solution}{iq:nw:disaster-recovery-and-regulatory-requirements:4}
L'annulation à la déconnexion (cancel on disconnect) et les kill switches des bourses arrêtent les ordres au repos~; les positions sont connues par les drop copies et le membre compensateur~;
les couvertures passent par le site de reprise ou par un courtier au téléphone~; puis le runbook relance le trading.
\iqlookfor{Exposition d'abord~; drop copies~; solutions de repli~; runbook.}
\end{solution}

\begin{solution}{iq:nw:disaster-recovery-and-regulatory-requirements:5}
À partir de la \omterm{def:nw:disaster-recovery-and-regulatory-requirements:resilience}{tolérance d'impact}~: simuler ou mesurer le délai de reprise de chaque dispositif par rapport à elle, chiffrer chacun et
comparer l'écart avec le coût attendu de l'indisponibilité supplémentaire.
\iqlookfor{Tolérance d'abord~; temps mesurés~; coût face à la perte attendue.}
\end{solution}

\begin{solution}{iq:nw:disaster-recovery-and-regulatory-requirements:6}
\omterm{def:nw:disaster-recovery-and-regulatory-requirements:resilience}{Services métier importants} et tolérances~; un site de reprise hors des risques du principal, à proximité des sites de reprise des plateformes~;
runbooks avec des durées mesurées~; arrêt ordonné et gestion des ordres en cours par plateforme~; tests annuels avec les bourses et le secteur~;
registres conservés comme preuves.
\iqlookfor{Règles rattachées aux dispositifs~; tests~; preuves.}
\end{solution}
